% Notas de aula — Capítulo 21: Processos Climáticos Naturais
% Curso: Meteorologia Prática (baseado em R. Stull, Practical Meteorology, CC BY-NC-SA 4.0)
% Caixas verdes = conteúdo literal do quadro; linhas verdes = encenação.
\documentclass[11pt]{article}
\usepackage{../comum/preambulo}
\graphicspath{{../figuras/cap21/}}

\title{Capítulo 21 --- Processos Climáticos Naturais\\
{\large Notas de aula (roteiro de quadro-negro)}}
\author{Curso de Meteorologia Prática}
\date{}

\begin{document}
\maketitle
\thispagestyle{fancy}

\noindent\textbf{Plano sugerido:} 2 aulas de 2\,h.
\emph{Aula 1:} \S\ref{sec:sistema}--\S\ref{sec:milan} (o sistema
climático, feedbacks, gelo-albedo, Milankovitch).
\emph{Aula 2:} \S\ref{sec:daisy}--\S\ref{sec:enso} (Daisyworld,
variabilidade interna, ENSO).

\section{Do tempo ao clima}\label{sec:sistema}

\begin{noquadro}[abertura]
\textbf{Cap.~21 --- Clima Natural}\\[2pt]
tempo $=$ a trajetória (imprevisível além de 2 semanas
\javimos{20}{caos})\\
clima $=$ a \emph{estatística} da trajetória — e ela é previsível\\[2pt]
personagens novos: \textbf{feedbacks}, \textbf{escalas lentas}
(oceano, gelo) e \textbf{forçantes externas} (Sol, órbita, vulcões)
\end{noquadro}

O balanço de energia do Cap.~2 \javimos{2}{$T_e = 255$\,K} volta como
protagonista — mas agora com o albedo, o vapor e as nuvens
\emph{respondendo} à própria temperatura: o sistema tem
retroalimentação, e retroalimentação muda tudo.

\section{A álgebra dos feedbacks}\label{sec:feed}

\quadro{Desenhar o diagrama de blocos do amplificador realimentado
(os alunos de eletrônica sorriem) e traduzir cada bloco para o
clima.}

\begin{formadif}[ganho com feedback (T1--T2)]
\[ \Delta T = \frac{\lambda_0\,\Delta F}{1 - f},\qquad
\lambda_0 = \frac{1}{4\varepsilon\sigma T^3} \simeq
0{,}3\ \mathrm{K/(W\,m^{-2})}. \]
Sem feedbacks, $2\times$CO$_2$ ($3{,}7$\,W/m²) dá $\sim$1,2\,K
(resposta de Planck). Com vapor ($f\sim+0{,}4$
\javimos{4}{Clausius--Clapeyron!}), gelo-albedo ($+0{,}1$--$0{,}2$) e
nuvens (incerto \javimos{6}{nuvens}): 2,5--4\,K.
\end{formadif}

\begin{noquadro}[o gelo-albedo em uma linha]
mais frio $\to$ mais gelo $\to$ mais albedo $\to$ mais frio\\
feedback POSITIVO $\Rightarrow$ múltiplos equilíbrios: para o MESMO
Sol, dois climas estáveis (Caso~1)\\[2pt]
a Terra bola-de-neve (~700\,Ma) visitou o outro ramo — e saiu pelo
CO$_2$ vulcânico (N1)
\end{noquadro}

A bifurcação do Caso~1 é o gráfico mais importante da aula: dois
ramos estáveis, um instável no meio, e histerese — sair do gelo exige
forçante muito além da que bastaria para não entrar.

\section{Milankovitch: o metrônomo orbital}\label{sec:milan}

\quadro{Os três ciclos com as mãos: excentricidade (abrir/fechar a
elipse, $\sim$100\,ka), obliquidade (inclinar o eixo, 41\,ka),
precessão (rodar o pião, 23/19\,ka).}

\begin{formalivro}[o gatilho das eras glaciais (T3; Caso~2)]
O que decide é a \textbf{insolação de verão em $\sim$65°N}
\javimos{2}{insolação por latitude}: se a neve do inverno sobrevive
ao verão, o gelo cresce. Precessão modula QUAL estação coincide com o
periélio; obliquidade, o contraste sazonal; excentricidade, a
amplitude da precessão.
\end{formalivro}

O mistério dos 100\,ka: a forçante quase não tem esse período, mas a
resposta tem — o gelo lento e não linear (Caso~1!) retifica o
metrônomo (N2): glaciação lenta, deglaciação abrupta. Vulcões e
variação solar completam as forçantes naturais: Pinatubo esfriou
$\sim$0,5\,K por 2 anos — os aerossóis \javimos{7}{aerossóis} em
escala planetária. Na escala mais lenta de todas ($10^7$--$10^8$
anos), a \textbf{tectônica} (Stull \S21.3) reposiciona continentes e
correntes: a Pangeia teve clima continental extremo, e a abertura da
passagem de Drake isolou (e congelou) a Antártida — geografia é
forçante climática.

\section{Daisyworld: regulação sem termostato}\label{sec:daisy}

\begin{formadif}[o planeta das margaridas (T4; Caso~3)]
Margaridas pretas (aquecem o planeta) prosperam com Sol fraco;
brancas (esfriam) com Sol forte. Sem intenção alguma, a cobertura
ajusta o albedo e a temperatura fica $\sim$ótima por um largo
intervalo de luminosidade: \textbf{feedback biológico negativo}.
A diversidade amplia a regulação (N3) — e o colapso nas bordas é
abrupto.
\end{formadif}

É o contraponto pedagógico do gelo-albedo: feedbacks positivos criam
saltos; negativos, platôs. O clima real tem os dois — e a pergunta
``qual domina?'' é a pergunta da sensibilidade climática.

\section{O clima presente e sua classificação}\label{sec:koppen}

\begin{noquadro}[Köppen em uma linha (Stull \S21.7)]
letras $=$ física: A tropical úmido (ZCIT) $\cdot$ B seco
(subsidência de Hadley) $\cdot$ C temperado $\cdot$ D continental
$\cdot$ E polar\\[2pt]
o mapa de Köppen é o mapa da circulação geral \javimos{11}{cinturões}
carimbado no chão.
\end{noquadro}

Cada zona de Köppen é diagnóstico das médias de $T$ e precipitação —
e portanto dos cinturões do Cap.~11: os B nos $30°$ de subsidência,
os A sob a ZCIT, os C onde as frentes \javimos{12}{frentes} visitam.
Deslocamentos das zonas são a régua de detecção de mudanças
\veremos{23}{mudança climática}. A hierarquia de modelos que fecha o
capítulo: EBM 0-D (Caso~1) $\to$ EBM meridional
\javimos{11}{climlab} $\to$ \textbf{GCMs} (Stull \S21.6) — os mesmos
modelos do Cap.~20 \javimos{20}{NWP} rodados por séculos, com oceano,
gelo e biosfera acoplados.

\section{Variabilidade interna: ENSO}\label{sec:enso}

\quadro{Desenhar o Pacífico equatorial: alísios
\javimos{11}{alísios}, piscina quente no oeste, ressurgência fria no
leste — e perguntar o que acontece se os alísios fraquejarem.}

\begin{formadif}[o oscilador de recarga (T5; Caso~4)]
TSM do leste ($T$) e conteúdo de calor ($h$):
\[ \dot T = aT + bh,\qquad \dot h = -cT - dh
\quad\Rightarrow\quad \text{oscilação de 3--5 anos}. \]
Acoplamento de Bjerknes ($a$): alísios fracos $\to$ leste quente
$\to$ alísios mais fracos — feedback positivo saturado pela descarga
lenta de calor. Ruído do \emph{tempo} torna os eventos irregulares.
\end{formadif}

\begin{noquadro}[por que ENSO importa aqui]
El Niño remodela as chuvas do planeta (Brasil: seca no N/NE, chuva no
S)\\
previsibilidade CLIMÁTICA de meses \javimos{20}{apesar do caos}:
o oceano lento é a memória\\
irmãos: PDO, AMO, NAO — o espectro contínuo da variabilidade interna
\end{noquadro}

\section{Síntese da aula}

\begin{noquadro}[mapa final --- canto do quadro]
\[ \Delta T = \frac{\lambda_0\Delta F}{1 - f} \qquad
\text{gelo-albedo: 2 equilíbrios} \qquad
Q_{65N}(\text{órbita}) \qquad
\dot T = aT + bh \]
feedbacks $+$ escalas lentas $+$ forçantes externas $=$ o clima que
varia sozinho — o palco onde o módulo de mudança climática (Cap.~23) entra
\end{noquadro}

\section{Exercícios complementares}\label{sec:exercicios}

Soluções (professor) em \texttt{cap21\_solucoes.ipynb}.

\subsection*{Teóricos}

\begin{enumerate}[label=\textbf{T\arabic*.},itemsep=4pt]
  \item Monte a álgebra de feedbacks (série geométrica) e classifique
        vapor d'água, gelo-albedo, Planck e Daisyworld em $f>0$ ou
        $f<0$.
  \item Calcule a resposta de Planck
        $\lambda_0 = 1/(4\varepsilon\sigma T^3)$ e o aquecimento sem
        feedbacks para $2\times$CO$_2$.
  \item Explique por que a insolação de verão em 65°N comanda as
        glaciações e associe cada ciclo orbital a seu período.
  \item Mostre que o equilíbrio interior do Daisyworld é estável
        (sinal do feedback via $dA/dT$) e discuta os colapsos de
        borda.
  \item Obtenha os autovalores do oscilador de recarga e a condição
        de oscilação; calcule o período para os parâmetros do
        Caso~4.
\end{enumerate}

\subsection*{Numéricos (no notebook)}

\begin{enumerate}[label=\textbf{N\arabic*.},itemsep=4pt]
  \item Simule o resgate vulcânico da bola de neve (histerese
        completa).
  \item Retifique Milankovitch com gelo lento $+$ limiar e recupere a
        periodicidade de $\sim$100\,ka.
  \item Compare o Daisyworld com uma e duas espécies: a diversidade
        amplia a regulação?
  \item Varra o acoplamento do ENSO e localize a bifurcação de Hopf;
        discuta o papel do ruído no regime subcrítico.
\end{enumerate}

\vspace{1em}
\noindent\rule{\linewidth}{0.4pt}\\
{\scriptsize Material derivado de R.~Stull, \emph{Practical Meteorology}
(CC BY-NC-SA 4.0); este documento herda a mesma licença.}

\end{document}
